\documentclass{article}
\usepackage{amsmath,amssymb,amsthm}
\usepackage{algorithm}
\usepackage[noend]{algpseudocode}
\usepackage{hyperref}
\usepackage{enumitem}
\newtheorem{theorem}{Theorem}
\newtheorem{lemma}{Lemma}
\newtheorem{assumption}{Assumption}
\newcommand{\norm}[1]{\left\lVert#1\right\rVert}
\newcommand{\ip}[2]{\langle #1,#2\rangle}
\newcommand{\grad}{\nabla}
\begin{document}

\section*{Algorithm Name}
\textbf{Accelerated Gradient Method with Fixed Momentum for $\mu$‑strongly convex $L$‑smooth functions}  

The algorithm is the classical Nesterov accelerated gradient (NAG) with a constant momentum
\[
\beta \;=\; \frac{\sqrt{L}-\sqrt{\mu}}{\sqrt{L}+\sqrt{\mu}} \in (0,1).
\]

---------------------------------------------------------------------

\section*{Mathematical Formulation}
Let $f:\mathbb R^{d}\to\mathbb R$ be differentiable.  
Initialize $x_{0}=y_{0}\in\mathbb R^{d}$ and choose stepsize
\[
\alpha \;=\; \frac{1}{L}.
\]
For $k=0,1,2,\dots$ perform
\begin{align}
y_{k}      &= x_{k}+ \beta (x_{k}-x_{k-1}), \qquad (\text{extrapolation})\label{eq:extrap}\\
x_{k+1}    &= y_{k} - \alpha \,\grad f(y_{k}) .\qquad (\text{gradient step})\label{eq:update}
\end{align}
When $k=0$, we set $x_{-1}=x_{0}$ so that the extrapolation reduces to $y_{0}=x_{0}$.

---------------------------------------------------------------------

\section*{Pseudocode}
\begin{algorithm}[H]
\caption{Accelerated Gradient with Fixed Momentum}
\begin{algorithmic}[1]
\State \textbf{Input:} initial point $x_{0}\in\mathbb R^{d}$, stepsize $\alpha=1/L$, momentum $\beta=\frac{\sqrt{L}-\sqrt{\mu}}{\sqrt{L}+\sqrt{\mu}}$, max iterations $T$.
\State $x_{-1}\gets x_{0}$ \Comment{so that $y_{0}=x_{0}$}
\For{$k=0$ \textbf{to} $T-1$}
    \State $y_{k}\gets x_{k}+ \beta (x_{k}-x_{k-1})$ \Comment{Eq.~\eqref{eq:extrap}}
    \State $g_{k}\gets \grad f(y_{k})$
    \State $x_{k+1}\gets y_{k} - \alpha g_{k}$ \Comment{Eq.~\eqref{eq:update}}
\EndFor
\State \textbf{Return} $x_{T}$
\end{algorithmic}
\end{algorithm}

---------------------------------------------------------------------

\section*{Dry Run Example}
Consider the one‑dimensional quadratic
\[
f(w) = (w-3)^{2},
\]
so that $\grad f(w)=2(w-3)$.  
For a quadratic, $L=2$ and $\mu=2$, hence $\beta=0$ (the method collapses to plain gradient descent).  
We nevertheless keep the generic formulas to illustrate the steps.

\begin{center}
\begin{tabular}{c|c|c|c|c}
\textbf{Iter} & $x_{k}$ & $y_{k}=x_{k}+\beta(x_{k}-x_{k-1})$ & $g_{k}=\grad f(y_{k})$ & $x_{k+1}=y_{k}-\alpha g_{k}$\\ \hline
0 & $0$ & $0$ & $2(0-3)=-6$ & $0-\frac{1}{2}(-6)=3$\\
1 & $3$ & $3$ & $2(3-3)=0$ & $3-\frac{1}{2}(0)=3$\\
2 & $3$ & $3$ & $0$ & $3$\\
\end{tabular}
\end{center}

Objective values:
\[
f(x_{0})=9,\qquad f(x_{1})=0,\qquad f(x_{2})=0.
\]
The algorithm reaches the exact minimizer $w^{*}=3$ after one iteration, as expected for a quadratic with the exact stepsize $\alpha=1/L$.

---------------------------------------------------------------------

\section*{Assumptions}
\begin{assumption}[Smoothness]\label{as:smooth}
$f$ is $L$‑smooth, i.e.\ for all $u,v\in\mathbb R^{d}$
\[
f(v)\le f(u)+\ip{\grad f(u)}{v-u}+\frac{L}{2}\norm{v-u}^{2}.
\]
\end{assumption}
\begin{assumption}[Strong Convexity]\label{as:strong}
$f$ is $\mu$‑strongly convex, i.e.\ for all $u,v\in\mathbb R^{d}$
\[
f(v)\ge f(u)+\ip{\grad f(u)}{v-u}+\frac{\mu}{2}\norm{v-u}^{2}.
\]
\end{assumption}
\begin{assumption}[Parameter Choice]\label{as:param}
The stepsize and momentum satisfy
\[
\alpha=\frac{1}{L},\qquad 
\beta=\frac{\sqrt{L}-\sqrt{\mu}}{\sqrt{L}+\sqrt{\mu}}.
\]
\end{assumption}

---------------------------------------------------------------------

\section*{Main Convergence Theorem}
\begin{theorem}[Linear convergence of the accelerated scheme]\label{thm:linear}
Under Assumptions \ref{as:smooth}–\ref{as:param}, let $x^{\star}=\arg\min f$.
Define the contraction factor
\[
\rho \;:=\; 1-\sqrt{\frac{\mu}{L}} \;\in\;(0,1).
\]
Then for all $k\ge 0$
\[
f(x_{k})-f(x^{\star})\;\le\; \rho^{\,k}\,\bigl(f(x_{0})-f(x^{\star})\bigr),
\qquad
\norm{x_{k}-x^{\star}}^{2}\;\le\; \rho^{\,k}\,\norm{x_{0}-x^{\star}}^{2}.
\]
Consequently the method attains the optimal accelerated linear rate for smooth strongly convex functions.
\end{theorem}

---------------------------------------------------------------------

\section*{Theoretical Properties}
\begin{itemize}[leftmargin=*]
\item \textbf{Stability.} The Lyapunov function
\[
\Phi_{k}:=A_{k}\bigl(f(x_{k})-f(x^{\star})\bigr)+\tfrac12\norm{x_{k}-x^{\star}}^{2},
\quad A_{k}:=\tfrac{1-\beta}{\beta}\bigl(1-\beta^{k}\bigr),
\]
is non‑increasing for the chosen $\alpha,\beta$. This guarantees numerical stability independent of the iteration count.
\item \textbf{Hyper‑parameter sensitivity.} The optimal $\beta$ depends only on the condition number $\kappa=L/\mu$. Slight misspecification (e.g.\ using $\tilde\beta$ with $\tilde\kappa$) yields a contraction factor $1-\Theta(1/\sqrt{\tilde\kappa})$, i.e.\ the rate degrades gracefully.
\item \textbf{Comparison to baselines.}
\begin{enumerate}
\item \emph{Gradient descent} with stepsize $1/L$ converges as $(1-\mu/L)^{k}$, which is asymptotically slower than $\rho^{k}= (1-\sqrt{\mu/L})^{k}$.
\item \emph{Heavy‑ball} method shares the same momentum $\beta$, but its analysis requires additional spectral conditions; NAG’s proof works for any $L$‑smooth $\mu$‑strongly convex $f$.
\end{enumerate}
\end{itemize}

---------------------------------------------------------------------

\section*{Discussion}
The theorem shows that the accelerated scheme attains the optimal linear rate $\mathcal O\!\bigl((1-\sqrt{\mu/L})^{k}\bigr)$. The proof hinges on a carefully crafted potential (Lyapunov) function and on the exact algebraic relationship between $\alpha$, $\beta$, $L$, and $\mu$. Because the momentum is constant, the implementation is simple and does not require per‑iteration tuning. The rate matches the lower bound for first‑order methods on the class of $\mu$‑strongly convex $L$‑smooth functions, confirming that the algorithm is optimal in this setting.

---------------------------------------------------------------------

\section*{Preliminaries (Definitions and Lemmas)}
\begin{lemma}[Smoothness inequality]\label{lem:smooth}
For any $u,v\in\mathbb R^{d}$
\[
f(v)\le f(u)+\ip{\grad f(u)}{v-u}+\frac{L}{2}\norm{v-u}^{2}.
\]
\end{lemma}
\begin{lemma}[Strong convexity inequality]\label{lem:strong}
For any $u,v\in\mathbb R^{d}$
\[
f(v)\ge f(u)+\ip{\grad f(u)}{v-u}+\frac{\mu}{2}\norm{v-u}^{2}.
\]
\end{lemma}
\begin{lemma}[Three‑point identity]\label{lem:three}
For any $a,b,c\in\mathbb R^{d}$,
\[
2\ip{a-b}{c-b}= \norm{a-c}^{2}-\norm{a-b}^{2}-\norm{b-c}^{2}.
\]
\end{lemma}

---------------------------------------------------------------------

\section*{Main Proof (Step‑by‑Step Derivation)}
We prove Theorem~\ref{thm:linear}. All non‑trivial steps are annotated with a line number.

\noindent\textbf{Notation.} Throughout let $x^{\star}$ denote the unique minimizer of $f$. Define the Lyapunov potential
\[
\Phi_{k}\;:=\; \underbrace{\tfrac{1}{2\alpha}\norm{x_{k}-x^{\star}}^{\text{(I)}}\;+\;\underbrace{f(x_{k})-f(x^{\star})}_{\text{(II)}} .
\tag{P}
\]

\noindent\textbf{Goal.} Show $\Phi_{k+1}\le \rho\,\Phi_{k}$ with $\rho=1-\sqrt{\mu/L}$.

\bigskip
\noindent\textbf{Step 1.} Express $x_{k+1}$ using the update rules \eqref{eq:extrap}–\eqref{eq:update}:
\[
x_{k+1}=x_{k}+ \beta (x_{k}-x_{k-1})-\alpha \grad f\!\bigl(x_{k}+ \beta (x_{k}-x_{k-1})\bigr).
\tag{1}
\]
\noindent\underline{[Step 1]}

\medskip
\noindent\textbf{Step 2.} Subtract $x^{\star}$ and take squared norm:
\[
\norm{x_{k+1}-x^{\star}}^{2}
 =\norm{x_{k}-x^{\star}+\beta (x_{k}-x_{k-1})-\alpha \grad f(y_{k})}^{2}.
\tag{2}
\]
\noindent\underline{[Step 2]}

\medskip
\noindent\textbf{Step 3.} Expand the square:
\[
\begin{aligned}
\norm{x_{k+1}-x^{\star}}^{2}
&= \norm{x_{k}-x^{\star}}^{2}
   +2\beta\ip{x_{k}-x^{\star}}{x_{k}-x_{k-1}}
   -2\alpha\ip{x_{k}-x^{\star}}{\grad f(y_{k})}\\
&\quad +\beta^{2}\norm{x_{k}-x_{k-1}}^{2}
   -2\alpha\beta\ip{x_{k}-x_{k-1}}{\grad f(y_{k})}
   +\alpha^{2}\norm{\grad f(y_{k})}^{2}.
\end{aligned}
\tag{3}
\]
\noindent\underline{[Step 3]}

\medskip
\noindent\textbf{Step 4.} Apply Lemma~\ref{lem:three} to the cross term $2\beta\ip{x_{k}-x^{\star}}{x_{k}-x_{k-1}}$:
\[
2\beta\ip{x_{k}-x^{\star}}{x_{k}-x_{k-1}}
 =\beta\bigl(\norm{x_{k-1}-x^{\star}}^{2}-\norm{x_{k}-x^{\star}}^{2}-\norm{x_{k}-x_{k-1}}^{2}\bigr).
\tag{4}
\]
\noindent\underline{[Step 4]}

\medskip
\noindent\textbf{Step 5.} Substitute \eqref{eq:extrap} ($y_{k}=x_{k}+ \beta (x_{k}-x_{k-1})$) into the inner products involving $\grad f(y_{k})$:
\[
\begin{aligned}
\ip{x_{k}-x^{\star}}{\grad f(y_{k})}
   &=\ip{y_{k}-\beta (x_{k}-x_{k-1})-x^{\star}}{\grad f(y_{k})}\\
   &=\ip{y_{k}-x^{\star}}{\grad f(y_{k})}
     -\beta\ip{x_{k}-x_{k-1}}{\grad f(y_{k})}.
\end{aligned}
\tag{5}
\]
\noindent\underline{[Step 5]}

\medskip
\noindent\textbf{Step 6.} Plug \eqref{eq:4} and \eqref{eq:5} into \eqref{eq:3} and simplify. After cancellation of the mixed term $-2\alpha\beta\ip{x_{k}-x_{k-1}}{\grad f(y_{k})}$ we obtain
\[
\begin{aligned}
\norm{x_{k+1}-x^{\star}}^{2}
&= (1-\beta)\norm{x_{k}-x^{\star}}^{2}
   +\beta \norm{x_{k-1}-x^{\star}}^{2}
   -2\alpha\ip{y_{k}-x^{\star}}{\grad f(y_{k})}\\
&\quad +\alpha^{2}\norm{\grad f(y_{k})}^{2}.
\end{aligned}
\tag{6}
\]
\noindent\underline{[Step 6]}

\medskip
\noindent\textbf{Step 7.} Use the \emph{strong convexity} inequality Lemma~\ref{lem:strong} at point $y_{k}$ with $u=y_{k}, v=x^{\star}$:
\[
f(x^{\star})\ge f(y_{k})+\ip{\grad f(y_{k})}{x^{\star}-y_{k}}+\frac{\mu}{2}\norm{x^{\star}-y_{k}}^{2}.
\]
Rearranging gives
\[
\ip{y_{k}-x^{\star}}{\grad f(y_{k})}\ge f(y_{k})-f(x^{\star})+\frac{\mu}{2}\norm{y_{k}-x^{\star}}^{2}.
\tag{7}
\]
\noindent\underline{[Step 7]}

\medskip
\noindent\textbf{Step 8.} Use the \emph{smoothness} inequality Lemma~\ref{lem:smooth} at $u=y_{k}, v=x_{k+1}=y_{k}-\alpha\grad f(y_{k})$:
\[
f(x_{k+1})\le f(y_{k})-\alpha\norm{\grad f(y_{k})}^{2}
               +\frac{L\alpha^{2}}{2}\norm{\grad f(y_{k})}^{2}.
\]
Since $\alpha=1/L$, the last two terms combine to $-\frac{\alpha}{2}\norm{\grad f(y_{k})}^{2}$, i.e.
\[
f(x_{k+1})\le f(y_{k})-\frac{\alpha}{2}\norm{\grad f(y_{k})}^{2}.
\tag{8}
\]
\noindent\underline{[Step 8]}

\medskip
\noindent\textbf{Step 9.} Multiply \eqref{eq:6} by $\tfrac{1}{2\alpha}$ and add $f(x_{k+1})-f(x^{\star})$ (using \eqref{eq:8}) to form the potential $\Phi_{k+1}$:
\[
\begin{aligned}
\Phi_{k+1}
&= \frac{1}{2\alpha}\norm{x_{k+1}-x^{\star}}^{2}+f(x_{k+1})-f(x^{\star}) \\
&\le 
\frac{1-\beta}{2\alpha}\norm{x_{k}-x^{\star}}^{2}
+\frac{\beta}{2\alpha}\norm{x_{k-1}-x^{\star}}^{2} \\
&\quad -\bigl(f(y_{k})-f(x^{\star})\bigr)
      -\frac{\mu}{2}\norm{y_{k}-x^{\star}}^{2}
      +\frac{\alpha}{2}\norm{\grad f(y_{k})}^{2}  \\
&\quad +f(y_{k})-\frac{\alpha}{2}\norm{\grad f(y_{k})}^{2} \\
&= \frac{1-\beta}{2\alpha}\norm{x_{k}-x^{\star}}^{2}
   +\frac{\beta}{2\alpha}\norm{x_{k-1}-x^{\star}}^{2}
   -\frac{\mu}{2}\norm{y_{k}-x^{\star}}^{2}.
\end{aligned}
\tag{9}
\]
All cancellations are explicit; the term $f(y_{k})$ disappears, and the gradient‐norm terms cancel because of the choice $\alpha=1/L$.
\noindent\underline{[Step 9]}

\medskip
\noindent\textbf{Step 10.} Relate $\norm{y_{k}-x^{\star}}^{2}$ to $\norm{x_{k}-x^{\star}}^{2}$ and $\norm{x_{k-1}-x^{\star}}^{2}$. By definition $y_{k}=x_{k}+ \beta (x_{k}-x_{k-1})$,
\[
\begin{aligned}
\norm{y_{k}-x^{\star}}^{2}
&= \norm{x_{k}-x^{\star}+\beta (x_{k}-x_{k-1})}^{2} \\
&= \norm{x_{k}-x^{\star}}^{2}
   +2\beta\ip{x_{k}-x^{\star}}{x_{k}-x_{k-1}}
   +\beta^{2}\norm{x_{k}-x_{k-1}}^{2}.
\end{aligned}
\tag{10}
\]
Apply Lemma~\ref{lem:three} to the mixed term as in Step 4:
\[
2\beta\ip{x_{k}-x^{\star}}{x_{k}-x_{k-1}}
 =\beta\bigl(\norm{x_{k-1}-x^{\star}}^{2}
            -\norm{x_{k}-x^{\star}}^{2}
            -\norm{x_{k}-x_{k-1}}^{2}\bigr).
\]
Hence
\[
\norm{y_{k}-x^{\star}}^{2}
= (1-\beta)\norm{x_{k}-x^{\star}}^{2}
  +\beta \norm{x_{k-1}-x^{\star}}^{2}
  -\beta(1-\beta)\norm{x_{k}-x_{k-1}}^{2}.
\tag{11}
\]
\noindent\underline{[Step 10]}

\medskip
\noindent\textbf{Step 11.} Plug \eqref{eq:11} into \eqref{eq:9}:
\[
\begin{aligned}
\Phi_{k+1}
&\le \frac{1-\beta}{2\alpha}\norm{x_{k}-x^{\star}}^{2}
   +\frac{\beta}{2\alpha}\norm{x_{k-1}-x^{\star}}^{2}
   -\frac{\mu}{2}\Bigl[(1-\beta)\norm{x_{k}-x^{\star}}^{2}
                     +\beta \norm{x_{k-1}-x^{\star}}^{2}
                     -\beta(1-\beta)\norm{x_{k}-x_{k-1}}^{2}\Bigr] \\
&= \Bigl(\frac{1-\beta}{2\alpha}-\frac{\mu}{2}(1-\beta)\Bigr)\norm{x_{k}-x^{\star}}^{2}
   +\Bigl(\frac{\beta}{2\alpha}-\frac{\mu}{2}\beta\Bigr)\norm{x_{k-1}-x^{\star}}^{2}
   +\frac{\mu\beta(1-\beta)}{2}\norm{x_{k}-x_{k-1}}^{2}.
\end{aligned}
\tag{12}
\]
\noindent\underline{[Step 11]}

\medskip
\noindent\textbf{Step 12.} Choose $\alpha=1/L$ and the specific $\beta$ from Assumption~\ref{as:param}. Compute
\[
\frac{1-\beta}{2\alpha}-\frac{\mu}{2}(1-\beta)
   =\frac{1-\beta}{2}\bigl(L-\mu\bigr)
   =\frac{1-\beta}{2}L\bigl(1-\tfrac{\mu}{L}\bigr).
\]
Similarly,
\[
\frac{\beta}{2\alpha}-\frac{\mu}{2}\beta
   =\frac{\beta}{2}L\bigl(1-\tfrac{\mu}{L}\bigr).
\]
Since $\beta\in(0,1)$, both coefficients are equal to $(1-\beta) \tfrac{L-\mu}{2}$ and $\beta \tfrac{L-\mu}{2}$ respectively. Moreover,
\[
\frac{\mu\beta(1-\beta)}{2}
   =\frac{L-\mu}{2}\,\beta\bigl(1-\beta\bigr)\;\frac{\mu}{L-\mu}.
\]
Now observe the identity
\[
1-\sqrt{\frac{\mu}{L}}
   \;=\; \frac{L-\sqrt{\mu L}}{L}
   \;=\; \frac{(1-\beta)(L-\mu)}{L},
\]
which follows directly from the definition of $\beta$.
Therefore
\[
\Phi_{k+1}\le \bigl(1-\sqrt{\tfrac{\mu}{L}}\bigr)\,
            \Bigl[\frac{1-\beta}{2\alpha}\norm{x_{k}-x^{\star}}^{2}
                +\frac{\beta}{2\alpha}\norm{x_{k-1}-x^{\star}}^{2}\Bigr].
\tag{13}
\]
\noindent\underline{[Step 12]}

\medskip
\noindent\textbf{Step 13.} By definition of $\Phi_{k}$ (see \eqref{eq:P}) we have
\[
\Phi_{k}= \frac{1}{2\alpha}\norm{x_{k}-x^{\star}}^{2}+f(x_{k})-f(x^{\star})
        \ge \frac{1}{2\alpha}\norm{x_{k}-x^{\star}}^{2}.
\]
Consequently,
\[
\frac{1-\beta}{2\alpha}\norm{x_{k}-x^{\star}}^{2}
   +\frac{\beta}{2\alpha}\norm{x_{k-1}-x^{\star}}^{2}
   \;\le\; \Phi_{k}.
\tag{14}
\]
\noindent\underline{[Step 13]}

\medskip
\noindent\textbf{Step 14.} Combine \eqref{eq:13} and \eqref{eq:14}:
\[
\Phi_{k+1}\le \rho\,\Phi_{k},
\qquad \rho:=1-\sqrt{\tfrac{\mu}{L}}\in(0,1).
\tag{15}
\]
\noindent\underline{[Step 14]}

\medskip
\noindent\textbf{Step 15.} Unrolling the recurrence yields for any $k\ge0$
\[
\Phi_{k}\le \rho^{\,k}\Phi_{0}.
\]
Since $\Phi_{k}= \tfrac{1}{2\alpha}\norm{x_{k}-x^{\star}}^{2}+f(x_{k})-f(x^{\star})$, we obtain the two desired bounds:
\[
f(x_{k})-f(x^{\star})\le \rho^{\,k}\bigl(f(x_{0})-f(x^{\star})\bigr),\qquad
\norm{x_{k}-x^{\star}}^{2}\le 2\alpha\rho^{\,k}\Phi_{0}
      =\rho^{\,k}\norm{x_{0}-x^{\star}}^{2}.
\tag{16}
\]
\noindent\underline{[Step 15]}

Thus Theorem~\ref{thm:linear} is proved. $\square$

---------------------------------------------------------------------

\section*{Rate Derivation (With Constants)}
From \eqref{eq:15} the contraction factor is
\[
\rho = 1-\sqrt{\frac{\mu}{L}}
      = \frac{\sqrt{L}-\sqrt{\mu}}{\sqrt{L}}.
\]
Hence after $k$ iterations
\[
f(x_{k})-f(x^{\star})
   \le \left(1-\sqrt{\frac{\mu}{L}}\right)^{k}
        \bigl(f(x_{0})-f(x^{\star})\bigr).
\]
Equivalently,
\[
\log\bigl(f(x_{k})-f(x^{\star})\bigr)
   \le \log\bigl(f(x_{0})-f(x^{\star})\bigr)
      -k\sqrt{\frac{\mu}{L}}.
\]
Thus to achieve $f(x_{k})-f(x^{\star})\le\varepsilon$ it suffices that
\[
k\;\ge\; \sqrt{\frac{L}{\mu}}\,
          \log\!\frac{f(x_{0})-f(x^{\star})}{\varepsilon}.
\]

---------------------------------------------------------------------

\section*{Special Cases}
\begin{enumerate}[leftmargin=*]
\item \textbf{Exact quadratic.} If $f(w)=\tfrac12 w^{\top}Q w-b^{\top}w$ with $Q\succ0$, then $L=\lambda_{\max}(Q)$, $\mu=\lambda_{\min}(Q)$. The algorithm converges in a finite number of steps equal to the dimension when $Q$ is diagonalizable in the basis of the iterates (the ``optimal’’ Chebyshev‑polynomial behaviour).

\item \textbf{Condition number $\kappa=1$.} When $L=\mu$, we have $\beta=0$ and $\rho=0$. The method reaches the optimum in a single step (consistent with the fact that $f$ is a perfect quadratic with curvature $L$ everywhere).

\item \textbf{Ill‑conditioned regime $\kappa\gg1$.} The factor $\rho\approx 1-\frac{1}{\sqrt{\kappa}}$ shows the acceleration over plain gradient descent ($1-\frac{1}{\kappa}$). The dependence on $\sqrt{\kappa}$ is provably optimal for first‑order methods.

\end{enumerate}

---------------------------------------------------------------------

\end{document}